\clearpage
\chapter*{강연 149의 정오표}
\addcontentsline{toc}{chapter}{강연 149의 정오표}
\setcounter{footnote}{0}
\src{298}

\noindent\textbf{1. 연접인 대수적 층에 대한 쌍대성 정리.}

\begin{center}
[부르바키 세미나, 제9권, 1956/57, 제149호, 25쪽.]
\end{center}

\noindent\textbf{14쪽에 관한 주.} --- 내가 1958년 국제수학자대회에서 한
강연에서 지적했듯이,\footnote{
Grothendieck (Alexander). --- The cohomology theory of abstract algebraic
varieties, \emph{Proceedings of the International Congress of Mathematicians}
[1958, Edinburgh], p. 103--118. --- Cambridge, at the University Press, 1960.}
여기서 제기된 문제들은 이제 완전히 해결되었다.

독자는 연접층의 쌍대성에 관한 더 자세한 내용을
\emph{loco citato}\footnotemark[1], 112--115쪽, EGA III 제2부(준비 중),
그리고 그때까지는 SGA 1962에서 찾을 수 있다. 더 체계적인 논의는 EGA의
뒤 이은 한 장(계획상 제IX장)에 실릴 것이다.
